\documentclass[11pt]{article}
\usepackage[T1]{fontenc}
\usepackage{lmodern}
\usepackage[margin=1in]{geometry}
\usepackage{amsmath,amssymb,amsthm,mathtools}
\usepackage{booktabs}
\usepackage{enumitem}
\usepackage{microtype}
\usepackage[hidelinks]{hyperref}

\newtheorem{theorem}{Theorem}[section]
\newtheorem{proposition}[theorem]{Proposition}
\newtheorem{lemma}[theorem]{Lemma}
\newtheorem{corollary}[theorem]{Corollary}
\theoremstyle{definition}
\newtheorem{definition}[theorem]{Definition}
\newtheorem{example}[theorem]{Example}
\theoremstyle{remark}
\newtheorem{remark}[theorem]{Remark}

\newcommand{\eps}{\varepsilon}
\newcommand{\D}{\mathcal D}
\newcommand{\im}{\operatorname{im}}
\newcommand{\cmp}{\operatorname{cmp}}
\newcommand{\cd}{\operatorname{cd}}
\newcommand{\mcd}{\operatorname{mcd}}
\newcommand{\ifrag}{\operatorname{ifrag}}
\newcommand{\tfrag}{\operatorname{tfrag}}
\newcommand{\Lock}{\operatorname{Lock}}
\newcommand{\Safe}{\operatorname{Safe}}
\newcommand{\Pareto}{\operatorname{Pareto}}
\newcommand{\blank}{\square}
\newcommand{\relto}{\mathrel{\rightsquigarrow}}

\title{Observation Is Not Enough:\\
Compression and Reconstruction under Finite-Monoid Typing}
\author{Takayuki Kuriyama\\Independent Researcher, Tokyo, Japan\\\texttt{growup.kuriyama@gmail.com}}
\date{}

\begin{document}
\maketitle

\begin{abstract}
Finite-monoid typing plays two different roles in positive-data grammar
learning.  It is a semantic compression device: a finite observer must retain
enough information to make guarded substitution safe.  It is also part of a
reconstruction architecture: once an observer is fixed, its fibers determine
which positive witnesses can be reused.  These two optimization problems need
not agree.

This note makes the separation precise by connecting finite-monoid compression
of principal syntactic-concept-lattice geometry with observer-sensitive
characteristic data.  First, observer refinement is shown to be monotone in the
opposite direction from the informal slogan that more information makes
learning easier: finer observers preserve safety but remove substitution
rules, split residual and transition cells, and can only increase minimum
locking-sample cost for the fixed reconstruction systems considered here.
Second, for the rigid CFG family $R_{m,q}$, every compression-optimal observer
has strictly larger characteristic-data cost than a slightly larger observer;
the exact gap is $4(m-1)(q-1)$.  In the regime $q=m^2$, the relative observer
overhead tends to one while the characteristic-data ratio tends to two.
Third, for the fan-out-two family $X_{k,r}$, the Clark--Wurm compression
sequence is exactly $1,2,2,\ldots$, yet for $X_{2,r}$ the optimal encoded
reconstruction cost is $2r(r+1)$.  Hence no function of the compression number
alone can bound reconstruction cost.  The same family also separates the full
Clark--Wurm tuple interface from the nonempty sentence-context interface used
by the MCFG learner.  Finally, the relational-morphism characterization of
minimum compression is lifted to a reconstruction separation: every selector
of every minimum-codomain separating relational morphism for $R_{m,q}$ remains
reconstruction-suboptimal.

The resulting picture is an observer-resource Pareto geometry in which the
finite-monoid compression number is only the first coordinate.  Minimal
information and minimal positive reconstruction are therefore distinct finite
resources.
\end{abstract}

\section{Introduction}

Distributional grammar learning reconstructs grammatical structure from
reusable contexts.  A finite monoid morphism
\[
  h:\Sigma^*\to M
\]
can be used as a finite typing bias: two fragments are candidates for
substitution only when their $h$-types agree.  In fixed-observer CFG and MCFG
learning this mechanism gives exact positive-data reconstruction for suitable
substitutable targets; see \cite{KuriyamaCFG,KuriyamaMCFG} and the earlier
substitutability literature \cite{ClarkEyraud2007,Yoshinaka2011}.

A separate question is how much finite information is actually necessary for
safe substitution.  In \cite{KuriyamaSCL}, finite-monoid observation is treated
as a compression problem on the principal layer of the syntactic concept
lattice (SCL): among all observers that separate every unsafe tuple conflict up
to arity $f$, minimize the image size.  For regular languages this minimum is
also characterized by finite-index congruences and by separating relational
morphisms on the pointed syntactic monoid.

The reconstruction problem begins only after this semantic compression problem
has been solved.  The observer may split one semantic residual class into many
typed fibers; higher-rank transitions may then fragment multiplicatively, and
positive data must expose enough incidences, arithmetic directions, and
ordered witnesses to reconnect the pieces.  These phenomena are quantified in
\cite{KuriyamaOFET}.

The organizing question of this note is therefore
\[
 \boxed{\text{Does the smallest safe observation also give the cheapest reconstruction?}}
\]
The answer is no in general, but the failure has two logically different
forms.

\begin{enumerate}[label=(\roman*),leftmargin=*]
\item \emph{Optimizer separation.}  For $R_{m,q}$, the observers minimizing
      image size and those minimizing characteristic data are disjoint.
\item \emph{Scalar insufficiency.}  For $X_{2,r}$, a two-state observer is
      simultaneously compression-optimal and reconstruction-optimal, but the
      optimal reconstruction cost still grows as $\Theta(r^2)$.  Thus the
      compression number alone does not control reconstruction cost.
\end{enumerate}

These two statements are the precise sense in which observation is not enough.
The note also records two structural principles.  Along a refinement chain,
finer observation can only remove legal substitution rules and increase the
minimum positive evidence required by the fixed reconstruction systems.  Thus
a strict image-size/data-cost tradeoff can occur only between refinement-
incomparable observers.  At the multiobjective level, the finite-monoid
compression number is the first-coordinate projection of a finite Pareto
frontier of observer resources.

\paragraph{Provenance.}
The exact $R_{m,q}$ fragmentation and characteristic-data formulas and the
exact $X_{k,r}$ slot-observer characteristic-data formula are established in
\cite{KuriyamaOFET}.  The SCL safety interface, congruence characterization,
relational-morphism characterization, and selector lemma are established in
\cite{KuriyamaSCL}.  The new content here is their exact connection: the two
interface compression numbers, the $X_{k,r}$ Clark--Wurm compression
calculation, the compression--reconstruction separation theorems, the
relational-selector separation, and the general refinement/Pareto principles
in the present form.

\section{Two tuple interfaces and two compression numbers}
\label{sec:interfaces}

For $L\subseteq\Sigma^*$ and a word $x\in\Sigma^*$, write
\[
 \D_L(x)=\{(u,v):uxv\in L\}.
\]
For an ordered $d$-tuple $\mathbf x=(x_1,\ldots,x_d)$, its tuple distribution
depends on the chosen class of tuple contexts.

\subsection{The Clark--Wurm interface}

A Clark--Wurm context of arity $d$ is
\[
 E=x_0\blank_1x_1\blank_2\cdots x_{d-1}\blank_dx_d,
 \qquad x_i\in\Sigma^*.
\]
Tuple components are also allowed to be empty.  Write
$\D^{(d),\mathrm{CW}}_L(\mathbf x)$ for the set of such contexts accepting
$\mathbf x$.  This is the fixed-order principal SCL interface used in
\cite{KuriyamaSCL}.

\begin{definition}[Clark--Wurm safety]
Let $h:\Sigma^*\to M$ be a finite monoid morphism.  Write
$\Safe_f^{\mathrm{CW}}(h;L)$ if for every $1\le d\le f$,
\[
 h^{(d)}(\mathbf x)=h^{(d)}(\mathbf y),\qquad
 \D_L^{(d),\mathrm{CW}}(\mathbf x)\cap
 \D_L^{(d),\mathrm{CW}}(\mathbf y)\ne\varnothing
\]
imply
\[
 \D_L^{(d),\mathrm{CW}}(\mathbf x)
 =
 \D_L^{(d),\mathrm{CW}}(\mathbf y).
\]
Define
\[
 \cmp_f^{\mathrm{CW}}(L)
 :=
 \min\{|\im h|:\Safe_f^{\mathrm{CW}}(h;L)\}.
\]
\end{definition}

By the principal-intent dictionary of \cite{KuriyamaSCL}, this is exactly the
finite-monoid compression number of the principal Clark--Wurm SCL geometry.
At arity one, equality of principal object concepts is equality modulo the
ordinary syntactic congruence, in agreement with Clark's SCL construction
\cite{Clark2013}.

\subsection{The positive sentence-context interface}

The MCFG reconstruction in \cite{KuriyamaMCFG,KuriyamaOFET} uses nonempty
ordered tuple components and sentence contexts
\[
 E=v_0\blank_1v_1\cdots v_{d-1}\blank_dv_d
\]
whose internal separators $v_1,\ldots,v_{d-1}$ are nonempty.  Let
$\D_L^{(d),+}$ denote this restricted distribution.

\begin{definition}[Positive-interface compression]
Write $\Safe_f^+(h;L)$ for the corresponding substitution condition on
nonempty tuples and positive internal separators, and put
\[
 \cmp_f^+(L)
 :=
 \min\{|\im h|:\Safe_f^+(h;L)\}.
\]
\end{definition}

This notation is introduced only to keep the learning interface distinct from
the full Clark--Wurm SCL interface.  For both compression numbers we use the
extended-natural convention that the minimum is $\infty$ if no safe finite
observer exists.

\begin{proposition}[Interface monotonicity]
\label{prop:interface-monotonicity}
For every language $L$ and $f\ge1$,
\[
 \cmp_f^+(L)\le \cmp_f^{\mathrm{CW}}(L).
\]
\end{proposition}

\begin{proof}
Every positive sentence context is a Clark--Wurm context, and every nonempty
tuple is a Clark--Wurm tuple.  Thus Clark--Wurm safety tests a superset of the
pairs tested by positive-interface safety.  Every Clark--Wurm-safe observer is
therefore positive-interface safe.
\end{proof}

The inequality can be strict; Section~\ref{sec:X} gives the exact family
$X_{k,r}$ with $\cmp_f^+=1$ and $\cmp_f^{\mathrm{CW}}=2$ for $f\ge2$.

\begin{lemma}[Fixed-length unary bridge]
\label{lem:fixed-length-bridge}
If $L\subseteq\Sigma^n$ for one fixed $n$, then empty and nonempty unary
fragments never form an overlapping Clark--Wurm pair.  Consequently the
empty-fragment convention does not change unary safety, and
\[
 \cmp_1^{\mathrm{CW}}(L)=\cmp_1^+(L).
\]
\end{lemma}

\begin{proof}
If one unary context $u\blank v$ accepted both $\eps$ and a nonempty word $x$,
then
\[
 |u|+|v|=n=|u|+|x|+|v|,
\]
which is impossible.  Hence all additional Clark--Wurm comparisons involving
$\eps$ have empty overlap and impose no unary safety constraint.
\end{proof}

\section{Observer refinement and reconstruction monotonicity}
\label{sec:refinement}

\begin{definition}[Observer refinement]
For finite monoid morphisms $h:\Sigma^*\to M$ and $g:\Sigma^*\to N$, write
\[
 h\preceq g
\]
when, after restricting codomains to their images, there exists a monoid
morphism $\pi:\im g\to\im h$ with
\[
 h=\pi\circ g.
\]
Thus $g$ is finer and $h$ is coarser.
\end{definition}

\begin{proposition}[Elementary refinement monotonicity]
\label{prop:elementary-refinement}
If $h\preceq g$, then
\[
 |\im h|\le|\im g|.
\]
Moreover, for either tuple interface $\mathsf I\in\{\mathrm{CW},+\}$,
\[
 \Safe_f^{\mathsf I}(h;L)\Longrightarrow
 \Safe_f^{\mathsf I}(g;L).
\]
If $R$ is any semantic residual class and
$\operatorname{frag}_h(R):=|h(R)|$, then
\[
 \operatorname{frag}_h(R)\le \operatorname{frag}_g(R).
\]
For pure witnessed higher-rank transitions, the transition fragmentation of
\cite{KuriyamaOFET} is likewise nondecreasing under refinement.
\end{proposition}

\begin{proof}
The factor map $\pi:\im g\twoheadrightarrow\im h$ gives the image-size
inequality.  Equality of $g$-types implies equality of $h$-types, so every
comparison required under $g$ is among those already controlled under $h$.
The residual-fiber inequality follows because the $g$-partition refines the
$h$-partition.  The transition statement follows from the product law
\[
 \tfrag_h(\theta)=\prod_j \operatorname{frag}_h(R_j)
\]
for pure witnessed transitions \cite{KuriyamaOFET}.
\end{proof}

We next isolate the reconstruction property used by both fixed-observer
learners in \cite{KuriyamaOFET}.  Their observer-dependent rules are precisely
semantic substitution rules controlled by equality of observer type; start,
constant, structural decomposition, and composition rules are observer
independent.

\begin{theorem}[Reconstruction monotonicity]
\label{thm:reconstruction-monotonicity}
Let $\mathcal A_h$ be either the fixed-observer CFG reconstruction or the
fixed-observer occurrence-MCFG reconstruction of \cite{KuriyamaOFET}.  If
$h\preceq g$, then for every finite positive sample $K$,
\[
 L(\mathcal A_g(K))\subseteq L(\mathcal A_h(K)).
\]
If $L$ is safe under $h$ for the tuple interface used by $\mathcal A_h$,
let $\Lock_h(L)$ be the set of characteristic locking samples for
$\mathcal A_h$.  Then
\[
 \boxed{\Lock_g(L)\subseteq\Lock_h(L).}
\]
Consequently, for every nonnegative sample cost $c(C)$ independent of the
observer,
\[
 \min_{C\in\Lock_h(L)}c(C)
 \le
 \min_{C\in\Lock_g(L)}c(C).
\]
\end{theorem}

\begin{proof}
Every semantic substitution rule available under $g$ is also available under
$h$, because equality of $g$-type implies equality of $h$-type.  All other
rules are unchanged.  Hence the $g$-hypothesis is a subgrammar of the
$h$-hypothesis.

Now let $C\in\Lock_g(L)$ and let $C\subseteq K\subseteq L$ be finite.
Because $h$ is safe, the $h$-learner is sound on positive subsets of $L$.
Therefore
\[
 L=L(\mathcal A_g(K))
 \subseteq L(\mathcal A_h(K))
 \subseteq L,
\]
so equality holds and $C\in\Lock_h(L)$.  Minimizing any fixed nonnegative
cost over nested feasible sets gives the final inequality.
\end{proof}

\begin{corollary}[Least-safe-observer coincidence principle]
\label{cor:least-safe}
Fix the safety interface used by the reconstruction system.  Suppose its
safe observers for $L$ have a least element $h_\bot$ under refinement.  Then
$h_\bot$ minimizes image size within that feasible class and simultaneously
minimizes every reconstruction cost covered by
Theorem~\ref{thm:reconstruction-monotonicity}.
\end{corollary}

\begin{corollary}[Incomparability principle]
\label{cor:incomparability}
If safe observers $h,g$ satisfy
\[
 |\im h|<|\im g|
 \qquad\text{and}\qquad
 \cd_h(L)>\cd_g(L),
\]
then $h$ and $g$ are refinement-incomparable.
\end{corollary}

\begin{proof}
If $h\preceq g$, reconstruction monotonicity gives the opposite cost
inequality.  If $g\preceq h$, image-size monotonicity gives the opposite image
inequality.
\end{proof}

Thus a genuine larger-observer/cheaper-reconstruction effect cannot be a
phenomenon along one refinement chain.  It must reallocate finite algebraic
information in incomparable ways.

\section{Exact optimizer separation on the CFG family $R_{m,q}$}
\label{sec:R}

Fix $m,q\ge2$ and let
\[
 P_q=\{p_0,\ldots,p_{q-1}\},\qquad
 A_m=\{a_0,\ldots,a_{m-1}\},
\]
with fresh symbols $z_{s,i}$.  Define
\[
 R_{m,q}
 =
 \{p_sa_i a_j:s<q,\ i,j<m\}
 \cup
 \{p_sa_i z_{s,i}:s<q,\ i<m\}.
\]
Every target word has length three.  The second block consists of rigidity
witnesses.  In the main block all suffixes $a_i a_j$ belong to one unary
residual class $F$.  For a safe observer $h$, put
\[
 \kappa_h
 :=
 |\{h(a_i a_j):i,j<m\}|.
\]

The following exact facts are proved in \cite{KuriyamaOFET}.  We reproduce the
form needed here.

\begin{theorem}[Exact $R_{m,q}$ resource formulas, \cite{KuriyamaOFET}]
\label{thm:R-input}
For every $m,q\ge2$:
\begin{enumerate}[label=(\roman*),leftmargin=*]
\item
\[
 \cmp_1^+(R_{m,q})=qm.
\]
There is an observer $H^-$ of image size $qm$ with $\kappa_{H^-}=m$.
\item Every safe $h$ with $|\im h|=qm$ satisfies $\kappa_h\ge m$.
\item If $\kappa_h<m$, then $|\im h|\ge q(m+1)$, and an observer $H^+$
      attains $|\im H^+|=q(m+1)$ with $\kappa_{H^+}=1$.
\item For the fixed-observer CFG reconstruction,
\[
 \min\{|C|:C\in\Lock_h(R_{m,q})\}
 =qm+m^2+\kappa_h(q-1),
\]
and, under the encoded cost $\|C\|=\sum_{w\in C}(|w|+1)$,
\[
 \boxed{
 \cd_h(R_{m,q})
 =4\bigl(qm+m^2+\kappa_h(q-1)\bigr).
 }
\]
\end{enumerate}
\end{theorem}

The characteristic-data formula has a direct rectangular interpretation.  The
observer fibers inside $F$ form column blocks.  Within each block, a locking
sample must connect $q$ row vertices to all suffix vertices; minimum locking
samples are spanning trees of the corresponding complete bipartite graphs.
Thus splitting one critical observer fiber creates an exact positive-exposure
penalty.

Because $R_{m,q}\subseteq\Sigma^3$, Lemma~\ref{lem:fixed-length-bridge}
identifies its unary positive-interface and Clark--Wurm compression numbers.

\begin{theorem}[Observation--reconstruction separation]
\label{thm:main-separation}
For every $m,q\ge2$,
\[
 \boxed{\cmp_1^{\mathrm{CW}}(R_{m,q})=qm.}
\]
Among compression-optimal observers,
\[
 \boxed{
 \min_{\substack{h\ \mathrm{safe}\\|\im h|=qm}}
 \cd_h(R_{m,q})
 =4(m^2+2qm-m).
 }
\]
But over all safe observers,
\[
 \boxed{
 \min_{h\ \mathrm{safe}}\cd_h(R_{m,q})
 =4(m^2+qm+q-1),
 }
\]
and every observer attaining this latter value has image size at least
$q(m+1)>qm$.  Hence
\[
 \arg\min_h |\im h|
 \ \cap\
 \arg\min_h \cd_h(R_{m,q})
 =\varnothing.
\]
The exact characteristic-data separation is
\[
 \boxed{4(m-1)(q-1).}
\]
\end{theorem}

\begin{proof}
The equality $\cmp_1^{\mathrm{CW}}=qm$ follows from
Theorem~\ref{thm:R-input}(i) and Lemma~\ref{lem:fixed-length-bridge}.
If $|\im h|=qm$, part (ii) gives $\kappa_h\ge m$.  Substitution into the exact
formula in part (iv) gives
\[
 \cd_h\ge4(qm+m^2+m(q-1))
 =4(m^2+2qm-m),
\]
with equality at $H^-$.  For arbitrary safe $h$, $\kappa_h\ge1$, so the same
formula gives
\[
 \cd_h\ge4(qm+m^2+q-1),
\]
with equality at $H^+$.  Part (iii) shows that every observer with
$\kappa_h<m$, in particular every observer attaining the global data minimum,
has at least $q(m+1)$ image values.  Subtracting the two displayed costs gives
$4(m-1)(q-1)$.
\end{proof}

\begin{corollary}[Constant overhead, unbounded additive saving]
For $q=2$, increasing observer size from $2m$ to $2m+2$ reduces the optimal
encoded characteristic-data cost by
\[
 4(m-1).
\]
\end{corollary}

\begin{corollary}[Vanishing relative overhead, factor-two data gain]
\label{cor:factor-two}
Let $q=m^2$ and $L_m=R_{m,m^2}$.  A reconstruction-optimal observer has
\[
 \frac{|\im H^+|}{\cmp_1^{\mathrm{CW}}(L_m)}
 =1+\frac1m\longrightarrow1,
\]
while
\[
 \frac{
 \min_{|\im h|=\cmp_1^{\mathrm{CW}}(L_m)}\cd_h(L_m)
 }{
 \min_h\cd_h(L_m)
 }
 \longrightarrow2.
\]
\end{corollary}

\begin{proof}
Here $\cmp_1^{\mathrm{CW}}=qm=m^3$ and
$|\im H^+|=q(m+1)=m^3+m^2$.  The two characteristic-data values are
\[
 4(2m^3+m^2-m)
 \quad\text{and}\quad
 4(m^3+2m^2-1),
\]
whose ratio tends to $2$.
\end{proof}

\section{The fan-out-two family $X_{k,r}$}
\label{sec:X}

Fix $k,r\ge2$.  For $1\le j\le r$ and $1\le i\le k$, let $a_{j,i}$ and
$b_{j,i}$ be distinct terminals, and set
\[
 w_{\mathbf i}
 =a_{1,i_1}\cdots a_{r,i_r}
  b_{1,i_1}\cdots b_{r,i_r},
 \qquad
 \mathbf i=(i_1,\ldots,i_r)\in[k]^r,
\]
\[
 X_{k,r}=\{w_{\mathbf i}:\mathbf i\in[k]^r\}.
\]
Every word has length $2r$.  Let
\[
 S_r=A_1\cdots A_rB_1\cdots B_r
\]
be the skeleton obtained by erasing indices:
$\sigma(a_{j,i})=A_j$ and $\sigma(b_{j,i})=B_j$.  Every symbol in $S_r$ occurs
exactly once.

\subsection{Exact Clark--Wurm compression}

Let $E_2=\{1,z\}$ be the two-element monoid with $z^2=z$, and define
\[
 e:\Sigma^*\to E_2,
 \qquad
 e(w)=
 \begin{cases}
 1,&w=\eps,\\
 z,&w\ne\eps.
 \end{cases}
\]
Thus $e$ records only emptiness.

\begin{lemma}[Skeleton-block rigidity for Clark--Wurm contexts]
\label{lem:skeleton-rigidity}
Suppose two Clark--Wurm tuples $\mathbf x,\mathbf y$ have the same componentwise
emptiness pattern and share an accepting context in $X_{k,r}$.  Delete every
coordinate that is empty in both tuples and concatenate the fixed context
pieces on its two sides.  Partition the remaining holes into maximal blocks in
which consecutive holes are separated by the empty word.  Then:
\begin{enumerate}[label=(\roman*),leftmargin=*]
\item for each block, the concatenation of the component skeletons supplied by
      $\mathbf x$ and by $\mathbf y$ is the same interval of $S_r$;
\item that interval is uniquely located in $S_r$; and
\item in every other accepting context for $\mathbf x$, exactly the same holes
      of the block remain separated by the empty word and occupy that same
      interval.
\end{enumerate}
\end{lemma}

\begin{proof}
After the common empty coordinates are removed, every remaining component is
nonempty.  The skeleton of any such component is a nonempty contiguous
substring of
\[
 S_r=A_1\cdots A_rB_1\cdots B_r.
\]
Because every symbol of $S_r$ occurs exactly once, every nonempty substring of
$S_r$ has at most one location.  Hence, in any accepting filling, each literal
component of $\mathbf x$ occupies a uniquely determined interval of $S_r$.
The same is true of every nonempty fixed context piece.

Consider a maximal block of holes whose intervening fixed pieces are empty.
In the shared accepting context, the block occupies the interval between the
uniquely located fixed material immediately to its left and right (with the
beginning or end of $S_r$ used when the block is exterior).  Therefore the
concatenation of the component skeletons in the block is uniquely determined.
The same context accepts $\mathbf y$, so its block must fill the same interval,
proving (i)--(ii).

Now let $D$ be any other context accepting $\mathbf x$.  Each component of
$\mathbf x$ must again occur at its unique skeleton location.  Two consecutive
components that were separated by the empty word in the shared context occupy
adjacent intervals of $S_r$.  A nonempty separator between them in $D$ would
insert at least one skeleton symbol between two adjacent uniquely located
intervals, which is impossible.  Thus all separators internal to the block
remain empty in $D$, and the block occupies the same union of intervals.  This
proves (iii).
\end{proof}

\begin{theorem}[Exact Clark--Wurm compression of $X_{k,r}$]
\label{thm:X-cw-compression}
For every $k,r\ge2$,
\[
 \boxed{\cmp_1^{\mathrm{CW}}(X_{k,r})=1,}
\]
and for every $f\ge2$,
\[
 \boxed{\cmp_f^{\mathrm{CW}}(X_{k,r})=2.}
\]
\end{theorem}

\begin{proof}
We first prove that the emptiness observer $e$ is Clark--Wurm safe at every
arity.  Let $\mathbf x,\mathbf y$ have equal componentwise $e$-type and share
an accepting context.  Normalize common empty coordinates as in
Lemma~\ref{lem:skeleton-rigidity}.  Consider any other context $D$ accepting
$\mathbf x$.  By the lemma, every empty-separated run occupies the same
skeleton interval under $D$, and the same run of $\mathbf y$ has the same
concatenated skeleton.  Hence replacing $\mathbf x$ by $\mathbf y$ preserves
the entire skeleton $S_r$.

It remains to check indices.  Let $U\subseteq\{1,\ldots,2r\}$ be the
union of skeleton positions occupied by the tuple material.  The block lemma
shows that $U$ is the same for $\mathbf x$ and $\mathbf y$, and that every
accepting context for $\mathbf x$, including $D$, places its tuple material on
that same set $U$.

Fix a slot $j$ and consider its paired positions $A_j,B_j$.  If neither
position lies in $U$, both terminals are fixed by $D$, and acceptance of
$D[\mathbf x]$ already supplies a matching pair.  If exactly one position lies
in $U$, let its index in $\mathbf x$ be $i$.  In the original shared context,
the other position is fixed; acceptance of both fillings forces the inside
terminal of $\mathbf y$ to have the same index $i$.  Since $D[\mathbf x]$ is
accepted, the fixed counterpart supplied by $D$ also has index $i$, so the
pair remains matched after replacing $\mathbf x$ by $\mathbf y$.  Finally, if
both positions lie in $U$, acceptance of the original shared context filled by
$\mathbf y$ forces the two terminals supplied by $\mathbf y$ to have one and
the same index.  Thus every slot is consistent in $D[\mathbf y]$, so
$D[\mathbf y]\in X_{k,r}$.  Symmetry gives equality of tuple distributions.
Hence $e$ is safe and $\cmp_f^{\mathrm{CW}}\le2$ for every $f$.

At arity one, the trivial observer is safe.  Empty and nonempty fragments
cannot overlap by Lemma~\ref{lem:fixed-length-bridge}.  If two nonempty
fragments share an accepting context, their unique skeleton substrings occupy
the same interval.  The same slot-by-slot argument above shows equality of
unary distributions.  Hence $\cmp_1^{\mathrm{CW}}=1$.

For the arity-two lower bound, write the all-index-$1$ word as
\[
 s_1s_2\cdots s_{2r}.
\]
Let
\[
 \mathbf x=(\eps,s_1),\qquad
 \mathbf y=(s_1,\eps).
\]
They share the accepting context
\[
 E=\blank_1\blank_2s_2\cdots s_{2r}.
\]
But for
\[
 D=\blank_1s_2\cdots s_{2r-1}\blank_2s_{2r}
\]
we have $D[\mathbf y]\in X_{k,r}$ while
$D[\mathbf x]\notin X_{k,r}$.  Thus the trivial observer is not arity-two
safe.  Therefore $\cmp_f^{\mathrm{CW}}\ge2$ for every $f\ge2$, proving the
claim.
\end{proof}

\subsection{The positive interface collapses to one state}

The same family behaves differently under the nonempty sentence-context
interface.

\begin{theorem}[Interface separation on $X_{k,r}$]
\label{thm:X-interface}
For every $k,r\ge2$ and every $f\ge1$,
\[
 \boxed{\cmp_f^+(X_{k,r})=1.}
\]
Consequently, for every $f\ge2$,
\[
 \boxed{
 \cmp_f^+(X_{k,r})
 <
 \cmp_f^{\mathrm{CW}}(X_{k,r}).
 }
\]
\end{theorem}

\begin{proof}
Take the trivial observer.  If $d>2r$, there is no accepting $d$-tuple of
nonempty components, so safety is vacuous.  Let $1\le d\le2r$ and suppose two
nonempty $d$-tuples share an admissible sentence context.  For $d\ge2$, every
internal separator is nonempty.  Together with the unique occurrence of every
symbol of $S_r$, these separators anchor each component to the same skeleton
interval on both sides of the comparison.  For $d=1$, the same conclusion
holds from the two exterior context pieces and the endpoints of $S_r$; in the
root case both tuples are complete target words and have the unique root
context.

Now consider a paired slot $(A_j,B_j)$.  If at least one of its two positions
lies outside tuple material, the shared accepting context fixes the outside
terminal and therefore forces the inside index to be identical in the two
tuples.  If both positions lie inside tuple material, acceptance forces the two
terminals of that slot to carry the same index within each tuple, although the
common index may change between the tuples.  Every other accepting sentence
context for the first tuple must place the same component intervals at the same
unique skeleton locations, so the same slot-by-slot alternatives apply.
Replacing the first tuple by the second therefore preserves membership in
$X_{k,r}$.  Symmetry gives equality of positive-interface distributions, and
the one-element observer is safe.
\end{proof}

This strict inequality is why the two tuple interfaces should not be silently
identified: the Clark--Wurm SCL invariant and the MCFG learner's admissible
comparison bias answer slightly different optimization questions.

\subsection{Coincidence of the optimal observer, but not of the resources}

Let $h_{\mathrm{slot}}$ be the slot observer of \cite{KuriyamaOFET}, mapping
$a_{j,i}$ to $A_j$ and $b_{j,i}$ to $B_j$ in a finite truncation monoid.  In
this section $\mcd_h$ uses the OFET encoding
\[
 \|K\|_{\mathrm{OFET}}:=\sum_{w\in K}\max(1,|w|).
\]
This is deliberately distinguished from the current positive-sample-size
notation in \cite{KuriyamaMCFG}, which uses $\sum_{w\in K}(|w|+1)$.  The
cardinality statements below are independent of that bookkeeping convention.
Since every word of $X_{k,r}$ has length $2r$, the OFET cost of a sample is
exactly $2r$ times its cardinality.  The fixed-observer MCFG reconstruction of
\cite{KuriyamaOFET} satisfies
\[
 \boxed{
 \mcd_{h_{\mathrm{slot}}}(X_{k,r})
 =2r[1+r(k-1)],
 }
\]
and every characteristic sample has at least $1+r(k-1)$ words.

\begin{proposition}[Equality of the relevant learner rules]
\label{prop:X-learner-equality}
For every finite $K\subseteq X_{k,r}$, the semantic unary rules generated by
the trivial observer, by the emptiness observer $e$, and by
$h_{\mathrm{slot}}$ coincide.  Hence the resulting hypotheses coincide on
$X_{k,r}$:
\[
 \widehat G_{\mathrm{triv}}(K)
 =\widehat G_e(K)
 =\widehat G_{h_{\mathrm{slot}}}(K).
\]
\end{proposition}

\begin{proof}
The learner observes only nonempty tuple components, so $e$ is constant on all
observed components and gives the same equality test as the trivial observer.
If two observed tuples share an admissible context, the proof of
Theorem~\ref{thm:X-interface} shows that corresponding components occupy the
same skeleton intervals.  Hence they have equal slot type automatically.
Thus every shared-context pair admitted by the trivial observer is also
admitted by $h_{\mathrm{slot}}$, while the converse is immediate.  All
non-unary learner rules are observer independent.
\end{proof}

\begin{theorem}[Compression--reconstruction coincidence on $X_{k,r}$]
\label{thm:X-coincidence}
Among Clark--Wurm-safe observers through arity two, the two-state emptiness
observer $e$ is compression-optimal and simultaneously attains the minimum
MCFG characteristic-data cost:
\[
 |\im e|=\cmp_2^{\mathrm{CW}}(X_{k,r})=2,
\]
\[
 \boxed{
 \min_{h:\Safe_2^{\mathrm{CW}}(h;X_{k,r})}
 \mcd_h(X_{k,r})
 =2r[1+r(k-1)].
 }
\]
\end{theorem}

\begin{proof}
Compression optimality is Theorem~\ref{thm:X-cw-compression}.  By
Proposition~\ref{prop:X-learner-equality}, $e$ has exactly the slot-observer
hypothesis on every sample, so it attains the displayed OFET cost.  No safe
observer can do better: the trivial observer is a coarsening of every observer,
and it induces the same hypothesis as $e$ and $h_{\mathrm{slot}}$ on this
family.  Theorem~\ref{thm:reconstruction-monotonicity} therefore makes this
cost globally minimal.
\end{proof}

\begin{theorem}[Fixed compression, unbounded reconstruction]
\label{thm:fixed-compression-unbounded}
There is an infinite family of finite regular languages $L_r$ such that
\[
 \cmp_2^{\mathrm{CW}}(L_r)=2
\]
for every $r$, while the minimum characteristic-sample cardinality tends to
infinity and the minimum encoded reconstruction cost is quadratic.  Explicitly,
for
\[
 L_r=X_{2,r},
\]
one has
\[
 \boxed{
 \min_{h:\Safe_2^{\mathrm{CW}}(h;L_r)}
 \ \min_{C\in\Lock_h(L_r)} |C|
 =r+1,
 }
\]
and, under the OFET encoding just fixed,
\[
 \boxed{
 \min_{h:\Safe_2^{\mathrm{CW}}(h;L_r)}
 \mcd_h(L_r)
 =2r(r+1)=\Theta(r^2).
 }
\]
Consequently there is no function $F:\mathbb N\to\mathbb N$ satisfying
\[
 \min_{h:\Safe_2^{\mathrm{CW}}(h;L)}\mcd_h(L)
 \le F(\cmp_2^{\mathrm{CW}}(L))
\]
for every target $L$ in any reconstruction class containing all languages
$X_{2,r}$.
\end{theorem}

\begin{proof}
Apply Theorems~\ref{thm:X-cw-compression} and
\ref{thm:X-coincidence} with $k=2$.  The OFET lower bound gives at least
$r+1$ sample words, and its explicit locking sample attains this cardinality.
Because every target word has length $2r$, this is equivalent to the encoded
minimum $2r(r+1)$.  If such a function $F$ existed, then
$2r(r+1)\le F(2)$ for all $r$, a contradiction.
\end{proof}

\begin{remark}[What the unboundedness does and does not say]
The alphabets, target-word lengths, and witnessed rule ranks grow with $r$.
The theorem therefore rules out a bound by the compression number alone; it
does not rule out bounds that also depend on alphabet size, word length, rule
rank, or other reconstruction parameters.
\end{remark}

\begin{remark}[Finite exhaustive sanity check]
The ancillary script \texttt{xkr\_cw\_small\_audit.py} exhaustively enumerates
all Clark--Wurm and positive-interface factorizations for
$(k,r)=(2,2),(2,3),(3,2)$ over several arities.  It verifies the two-state
Clark--Wurm upper bound, the one-state positive-interface upper bound, and the
arity-two Clark--Wurm lower bound in those finite cases.  This computation is
only a sanity check; the all-arity result is the proof above.
\end{remark}

Thus $X_{k,r}$ gives the opposite phenomenon from $R_{m,q}$: here a
compression-optimal observer is reconstruction-optimal, yet the numerical
compression invariant still carries too little information to predict the
amount of positive evidence required.

\section{Relational morphisms optimize feasibility, not reconstruction}
\label{sec:relational}

Assume $L$ is regular, with syntactic morphism
\[
 \eta:\Sigma^*\twoheadrightarrow T_L
\]
and accepting subset $P\subseteq T_L$.  The relational-morphism theorem of
\cite{KuriyamaSCL} characterizes Clark--Wurm compression by
\[
 \cmp_f^{\mathrm{CW}}(L)
 =
 \min\{|M|:\rho:T_L\relto M\text{ is $f$-separating}\}.
\]
Here separation means that every unsafe syntactic tuple pair has a coordinate
whose relational fibers are disjoint.

The same paper proves a free-monoid selector lemma.  For each letter $a$, pick
one value in $\rho(\eta(a))$ and extend freely to a morphism
$h_{\rho,s}:\Sigma^*\to M$.  Then
\[
 \rho_{h_{\rho,s}}(t)\subseteq\rho(t)
\]
for every $t\in T_L$, so a separating relational morphism yields a safe
functional observer.

\begin{theorem}[Relational-morphism reconstruction separation]
\label{thm:relational-separation}
Let $L=R_{m,q}$ and let
\[
 \rho:T_L\relto M
\]
be any minimum-codomain unary separating relational morphism.  Then
$|M|=qm$.  Every free-monoid selector $h_{\rho,s}$ is compression-optimal:
\[
 |\im h_{\rho,s}|=qm.
\]
Nevertheless every such selector is reconstruction-suboptimal, and
\[
 \boxed{
 \cd_{h_{\rho,s}}(R_{m,q})
 -
 \min_h\cd_h(R_{m,q})
 \ge4(m-1)(q-1).
 }
\]
\end{theorem}

\begin{proof}
By the relational characterization and
Theorem~\ref{thm:main-separation}, $|M|=\cmp_1^{\mathrm{CW}}(R_{m,q})=qm$.
The selector lemma gives a safe observer with
$|\im h_{\rho,s}|\le |M|$.  A strict inequality would contradict minimality of
$\cmp_1^{\mathrm{CW}}$, so equality holds.  The compression-optimal lower bound
from Theorem~\ref{thm:main-separation} then applies to every selector, while the
global reconstruction optimum is smaller by exactly $4(m-1)(q-1)$.
\end{proof}

Thus relational morphisms solve the semantic feasibility problem exactly, but
a minimum relational codomain does not solve the positive reconstruction
problem.  For reconstruction one must descend from a multivalued relational
witness to the behavior of its functional selectors and their induced fibers.

\section{Observer-resource Pareto geometry}
\label{sec:pareto}

Fix one safety interface $\mathsf I$ (for example $\mathrm{CW}$ or $+$).
Let $F_1,\ldots,F_s$ be finite-valued observer resources, for example residual
fragmentation, transition fragmentation, minimum characteristic-sample
cardinality, or encoded characteristic-data cost.  Define
\[
 \mathcal R_L^{\mathsf I}(h)
 =
 \bigl(
 |\im h|,F_1(h),\ldots,F_s(h)
 \bigr)
 \in\mathbb N^{s+1}
\]
for $\mathsf I$-safe observers $h$, and order these vectors componentwise.
Write $\Pareto_f^{\mathsf I}(L)$ for the set of realized componentwise-minimal
vectors.

\begin{theorem}[Observer-resource Pareto principle]
\label{thm:pareto}
If $L$ admits at least one $\mathsf I$-safe finite observer and all chosen
resource coordinates are finite natural numbers, then
$\Pareto_f^{\mathsf I}(L)$ is finite and nonempty.  Moreover,
\[
 \boxed{
 \cmp_f^{\mathsf I}(L)
 =
 \min\{r_1:(r_1,\ldots,r_{s+1})\in\Pareto_f^{\mathsf I}(L)\}.
 }
\]
If all coordinates are nondecreasing under observer refinement, then two
distinct Pareto points cannot be represented by comparable observers.
\end{theorem}

\begin{proof}
The set of realized $\mathsf I$-resource vectors is a nonempty subset of
$\mathbb N^{s+1}$.  By Dickson's lemma, its set of componentwise minimal
elements is finite.  Every realized vector dominates a minimal one, so a vector
with globally minimum first coordinate dominates a Pareto-minimal vector with
the same first coordinate; this gives the displayed equality.  Finally, if
$h\preceq g$ and every coordinate is refinement-monotone, then
$\mathcal R_L(h)\le\mathcal R_L(g)$, so comparable observers cannot represent
two distinct Pareto-minimal vectors.
\end{proof}

For $R_{m,q}$, \cite{KuriyamaOFET} computes the full three-coordinate frontier
for
\[
 (|\im h|,\ifrag_h,\cd_h)
\]
as exactly
\[
 P^-_{m,q}
 =
 \bigl(qm,m,4(m^2+2qm-m)\bigr)
\]
and
\[
 P^+_{m,q}
 =
 \bigl(q(m+1),1,4(m^2+qm+q-1)\bigr).
\]
The first point is compression-optimal; the second is reconstruction-optimal.

\begin{corollary}[Weighted phase transition]
\label{cor:weighted}
For positive weights $\alpha,\beta,\gamma$, define
\[
 W_{\alpha,\beta,\gamma}(h)
 =
 \alpha|\im h|
 +\beta\ifrag_h(R_{m,q})
 +\gamma\cd_h(R_{m,q}).
\]
The minimizing Pareto profile is $P^-_{m,q}$ precisely when
\[
 \alpha q>
 (m-1)\bigl[\beta+4\gamma(q-1)\bigr],
\]
and is $P^+_{m,q}$ precisely when the reverse strict inequality holds.  At
equality both profiles minimize the weighted objective.
\end{corollary}

\begin{proof}
Every positive weighted optimum is Pareto-minimal, so only the two displayed
profiles need be compared.  Their objective difference is
\[
 W(P^-)-W(P^+)
 =
 -\alpha q
 +(m-1)\bigl[\beta+4\gamma(q-1)\bigr].
\]
\end{proof}

This phase transition is a nontrivial weighted extension of SCL compression:
changing the relative price of global observation and positive reconstruction
moves the optimum discontinuously between two algebraically incomparable
observers.

\section{Conclusion}

Finite observation and positive reconstruction are distinct resources.  The
compression number
\[
 \cmp_f(L)
\]
answers a semantic question: how many finite observer values are needed to
separate all unsafe guarded substitutions?  Characteristic-data cost answers a
different question: once those distinctions are implemented by a particular
observer, how many positive witnesses are required to reconstruct the target
under the fixed learner?

The family $R_{m,q}$ shows that these two optimization problems can have
disjoint optimizers, with an exact gap and a nontrivial two-point Pareto
frontier.  The family $X_{2,r}$ shows something stronger in a different
direction: even when a compression-optimal observer is also reconstruction-
optimal, the compression number can remain fixed while the optimal
reconstruction cost grows without bound.  The relational-morphism
characterization does not eliminate this distinction: minimum relational
codomain controls semantic feasibility, whereas reconstruction depends on the
functional observer geometry induced after selection.

The combined principle is therefore
\[
 \boxed{\textbf{Observation Is Not Enough}.}
\]
A finite observer can make reconstruction possible without determining how
expensive that reconstruction will be.  In the resulting observer-resource
geometry, SCL compression is the first coordinate, not the whole problem.

\begin{thebibliography}{99}

\bibitem{ClarkEyraud2007}
A.~Clark and R.~Eyraud.
\newblock Polynomial identification in the limit of substitutable context-free languages.
\newblock \emph{Journal of Machine Learning Research}, 8:1725--1745, 2007.

\bibitem{Clark2013}
A.~Clark.
\newblock The syntactic concept lattice: another algebraic theory of the context-free languages?
\newblock \emph{Journal of Logic and Computation}, 25(5):1203--1229, 2015.
\newblock Published online 2013.

\bibitem{Yoshinaka2011}
R.~Yoshinaka.
\newblock Efficient learning of multiple context-free languages with multidimensional substitutability from positive data.
\newblock \emph{Theoretical Computer Science}, 412(19):1821--1831, 2011.

\bibitem{KuriyamaCFG}
T.~Kuriyama.
\newblock Distributional learning of context-free languages under fixed finite-monoid typing.
\newblock Revised manuscript, 2026. Public predecessor: arXiv:1409.6247v4.

\bibitem{KuriyamaMCFG}
T.~Kuriyama.
\newblock Positive-data learning of multiple context-free languages under fixed finite-monoid typing.
\newblock arXiv:2605.11644, 2026.

\bibitem{KuriyamaSCL}
T.~Kuriyama.
\newblock Finite-monoid compression in syntactic concept lattices: arity hierarchies and a pseudovariety trichotomy.
\newblock arXiv:2608.29639, 2026.

\bibitem{KuriyamaOFET}
T.~Kuriyama.
\newblock Observer--Fragmentation--Exposure Tradeoffs: From rectangular CFG exposure to ordered MCFG scheduling.
\newblock arXiv:2609.34560, 2026.

\end{thebibliography}

\end{document}
